%%%%%%%%%%%%%%%%%%%%%%%%%%%%%%%%%%%%%%%%%%%%%%%%%%%%%%%%%%%%%
% ==================== 5.2：问题二的模型建立与求解 ====================
% ✓ 自查：问题一输出已真实进入本问；新增条件已体现；比较标准明确；检验不过度展开。

\subsection{问题二的模型建立与求解}

\subsubsection{建模思路与模型选择}

问题二在问题一的基础上进一步要求【概括问题二的任务】，新增了【新增条件或研究要求】。
本问以问题一得到的【参数、特征或中间结果】作为【本问中的作用】，重点解决【核心数学问题】。
根据【新增条件和数据特点】，选择【主模型名称】对问题一的方法进行【扩展、修正或重新建模】，
并采用【必要的对比模型或基准方案】从【比较指标】方面评价改进效果。
如需额外处理数据，则说明【处理对象、方法及理由】，不单独设置无实际内容的数据预处理小节。

\subsubsection{模型的建立}

在问题一符号体系的基础上，新增【变量或参数】表示【含义】，其余符号与前文保持一致。
依据【理论、机理或题目规则】建立本问的【核心方程、目标函数或判别规则】，并加入【新增约束或条件】：

\begin{equation}
  \text{【问题二核心方程、目标函数或判别规则】},
  \label{eq:q2_core}
\end{equation}

\begin{equation}
  \text{【问题二新增约束、边界条件或适用条件】}.
  \label{eq:q2_constraint}
\end{equation}

\noindent\textbf{模型汇总：}
综合问题一的输入、本问新增条件和变量范围，得到问题二的完整模型

\begin{equation}
  \mathcal{M}_2:\quad
  \left\{
  \begin{aligned}
    &\text{【核心模型表达式】},\\
    &\text{【继承条件与新增约束】},\\
    &\text{【变量范围、初始条件或边界条件】}.
  \end{aligned}
  \right.
  \label{eq:q2_summary}
\end{equation}

\subsubsection{模型的求解}

\noindent\textbf{Step 1：}读取问题一输出的【中间结果】并结合【新增数据或条件】，构造问题二的求解输入。

\noindent\textbf{Step 2：}采用【算法或计算方法】求解【主模型】，记录【目标值、状态量或评价指标】。

\noindent\textbf{Step 3：}在相同输入和评价口径下求解【对比模型或基准方案】，根据【比较指标】确定最终结果并输出【传递给问题三的内容】。

求解使用【软件、求解器或程序语言】完成，关键参数和终止条件为【真实设置及依据】，完整程序见附录【编号】。

\subsubsection{求解结果与分析}

问题二得到【真实的核心结果】。由【图或表编号】可知，【主模型】在【比较指标】方面表现为【真实对比结果】，
说明新增的【条件、机制或改进措施】对【研究对象】产生了【实际影响】。
据此，问题二的答案为【直接结论】，并将【参数、方案或预测结果】用于问题三。

\subsubsection{模型检验与灵敏度分析}

围绕【本问主要结论】，采用【与模型类型匹配的检验方法】验证【拟合、排序、约束或机制】是否可靠。
指标【指标名称】及其真实结果为【结果】，依据【判定标准或基准】可得【检验结论】。
仅当【关键参数】存在现实波动且可能改变结论时，分析其在【有依据的范围】内变化对
【核心输出或最优策略】的影响，并报告【稳定区间、敏感方向或策略切换点】。
